\subsection*{一步一步建立直观}

\paragraph{1. 从一个问题出发选择要读取的信息}

把查询看成当前要找的特征，把键看成可供匹配的索引，把值看成真正需要传递的内容。右图直接调两个匹配分数的差，省去向量点积，只观察匹配之后的读取阶段。分数相等时，两个值具有相同权重，输出落在它们的中点。



\paragraph{2. 权重移动，输出也随之移动}

提升第二项分数后，它的指数权重增大，输出向第二个值移动。由于权重非负且总和为一，单次标量加权结果始终落在两个值之间。读数与图形共同显示这个凸组合约束；经过后续线性层或残差相加后，整体输出则不再受同一界限约束。



\paragraph{3. 温度改变选择的锐利程度}

固定分数差而降低温度，较大分数占据更大份额；提高温度则趋向平均。softmax 仍然是连续加权，不是简单地取最大值。非常小的温度会使数值敏感，程序还应使用稳定计算方法。这里的两个标量值只是便于检查机制的缩小模型。



\subsection*{定义、命题与证明}

\paragraph{定义：归一化指数权重}

给定实数分数 $s_j$ 与正温度 $T$，令 $a_j=e^{s_j/T}/\sum_k e^{s_k/T}$。各权重为正，且总和为一。

\paragraph{命题：一次注意力读取是凸组合}

若标量值都在区间 $[m,M]$ 中，则加权输出 $\sum_j a_jv_j$ 也在 $[m,M]$ 中。

\begin{proof}

逐项有 $a_jm\le a_jv_j\le a_jM$，因为各权重非负。求和后得到 $m\sum_j a_j\le\sum_j a_jv_j\le M\sum_j a_j$。

使用权重和为一便得到结论。本命题只约束这一次加权汇总，未约束接下来的线性输出变换、残差相加或整个网络的数值范围。

\end{proof}
